\documentclass{article}

\begin{document}
    \begin{center}
        \Large \textbf{Challenge 1: Shared State} \\ [6pt]
        \normalsize
        Autor: Maximilian Jakob Maag B.Sc. \\
        Version: 1.0
    \end{center}
    \section{Intro}
    Terraform ist besonders gut geeignet, um Infrastrutkur zu provisonieren.
    Aufgaben nach der Provisionierung, wie z.B. Updates führen dazu dass Terraform ein VPS zerstört und neu aufsetzt.
    Dabei kann es zu Datenverlust kommen. Anstatt Terraform diese Aufgabe erledigen zu lassen kann z.B. Ansible verwendet werden.
    \\ \\
    Ansible ist in der Lage sich mit einem VPS zu verbinden und ein vordefiniertes Skript auszuführen.
    Eine Installation, ein Backup, ein Rollback oder ein Update kann so durchgeführt werden.

    \section{Aufgabe}
    In dieser Challenge sollen Sie eine möglichst große Anzahl an VPS deployen.
    In der Praxis ist die Zahl an Servern viel zu groß, um diese von Hand zu warten oder zu installieren. \\ \\
    Das Bereitstellen der VPS inklusive DNS Record und Firewall übernimmt Terraform. Alles andere ist die Aufgabe von Ansible.

    \subsection{Aufgabe 1}
    Erstellen Sie ein Modul in Terraform für ein generisches VPS. Dieses muss ein Volume besitzen. Das Volume  muss gemounted werden.
    Zusetzlich deployen Sie eine Firewall und  weisen das VPS der Firewall zu. Öffnen Sie geeignete Ports.

    \subsection{Aufgabe 2}
    Erstellen Sie ein Ansible Playbook, welches das VPS updated und einen Webserver installiert eine einfache Hello World  webseite soll angezeigt werden, sobald sich ein Nutzer mit dem Webserver verbindet.

    \subsection{Aufgabe 3}
    Passen Sie ihr Terraform Modul an, das VPS soll nun ebenfalls so vorbereitet werden, dass Sie Ihr Ansible Playbook auf jeder Instanz des VPS gleichzeitig starten können.

    \subsection{Aufgabe 4}
    Testen Sie Ihre Lösung, mit 2, 4, 8, 16 und 32 Instanzen Ihres VPS.
    Entwickeln Sie zusätzlich ein kleines Shellskript welches Sie einmalig mit der Anzahl der Instanzen aufrufen. \\ 
    Führen Sie nach der Challenge terraform destroy aus.
\end{document}